\documentclass[../../main.tex]{subfiles}

\begin{document}

\section{$L^{p}(X,\mu)(1\leqslant p\leqslant\infty)$的完备性}\label{section:ch19-completeness}

令 $(X,\mathcal{M},\mu)$ 为测度空间. 定义 $\mathcal{F}$ 为 $X$ 上几乎处处有限的可测扩充实值函数全体. 由于 $X$ 上的可积函数在 $X$ 上是 a.e. 有限的, 若 $f$ 是 $X$ 上的可测函数且存在 $(0,\infty)$ 中的 $p$ 使得 $\int_{X}|f|^{p}d\mu<\infty$, 则 $f$ 属于 $\mathcal{F}$. 定义 $\mathcal{F}$ 中的两个函数 $f$ 和 $g$ 为等价的, 若在 $X$ 上 $f=g$ a.e., 记为 $f\cong g$. 这是一个等价关系, 即它是自反、对称与传递的. 因此它诱导了 $\mathcal{F}$ 的一个等价类的不交族的分解. 我们记该等价类族为 $\mathcal{F}/\cong$. $\mathcal{F}/\cong$ 上存在自然的线性结构. 给定两个等价类 $[f]$ 与 $[g]$ 以及实数 $\alpha$ 和 $\beta$, 我们定义线性组合 $\alpha[f]+\beta[g]$ 为属于 $\mathcal{F}$ 的在 $X_{0}$ 上取值 $\alpha f(x)+\beta g(x)$ 的函数的等价类, 其中 $X_{0}$ 是 $X$ 中使得 $f$ 和 $g$ 都有限的点的集合. 观察到等价类的线性组合是恰当定义的, 而与等价类的代表的选取无关. 该线性空间的零元是在 $X$ 上几乎处处消失的函数的等价类.

对 $1\leqslant p<\infty$, 令 $L^{p}(X,\mu)$ 为使得
\begin{align*}
\int_{X}|f|^{p}\,d\mu<\infty
\end{align*}
的等价类 $[f]$ 的全体. 这是恰当定义的, 由于若 $f\cong f_{1}$, 则在 $X$ 上 $|f|^{p}$ 可积当且仅当 $|f_{1}|^{p}$ 如此. 我们从不等式
\begin{align*}
\text{对所有 }a,b\in\mathbf{R},\quad|a+b|^{p}\leqslant 2^{p}[|a|^{p}+|b|^{p}]
\end{align*}
和积分比较判别法推出 $L^{p}(X,\mu)$ 是线性空间. 对 $L^{p}(X,\mu)$ 中的等价类 $[f]$, 我们定义 $\|[f]\|_{p}$ 为
\begin{align*}
\|[f]\|_{p}=\left[\int_{X}|f|^{p}\,d\mu\right]^{1/p}.
\end{align*}
这是恰当定义的. 显然 $\|[f]\|_{p}=0$ 当且仅当 $[f]=0$, 且对每个实数 $\alpha$, $\|\alpha[f]\|_{p}=|\alpha|\|[f]\|_{p}$.

我们称等价类 $[f]$ 为本性有界的, 若存在 $M\geqslant 0$, 使得在 $X$ 上, $|f|\leqslant M$ a.e. $M$ 称为 $[f]$ 的本性上界. 这也是恰当定义的, 即不依赖于等价类的代表的选取. 我们定义 $L^{\infty}(X,\mu)$ 为使得 $f$ 本性有界的等价类 $[f]$ 的全体. 则 $L^{\infty}(X,\mu)$ 也是 $\mathcal{F}/\cong$ 的一个线性子空间. 对 $[f]\in L^{\infty}(X,\mu)$, 定义 $\|[f]\|_{\infty}$ 为 $f$ 的本性上界的下确界. 这是恰当定义的. 容易看出 $\|[f]\|_{\infty}$ 是 $f$ 的最小本性上界. 此外, $\|[f]\|_{\infty}=0$ 当且仅当 $[f]=0$, 且对每个实数 $\alpha$, $\|\alpha[f]\|_{\infty}=|\alpha|\|[f]\|_{\infty}$. 我们从关于实数的三角不等式推出关于 $\|\cdot\|_{\infty}$ 的三角不等式. 因此 $\|\cdot\|_{\infty}$ 是一个范数.

为简单与方便起见, 我们将 $\mathcal{F}/\cong$ 中的等价类视为函数, 记为 $f$ 而非 $[f]$. 因此写 $f=g$ 意味着对几乎所有 $x\in X$, $f(x)=g(x)$.

回忆 $(1,\infty)$ 中的数 $p$ 的共轭 $q$ 由关系式 $1/p+1/q=1$ 定义, 我们也称 $1$ 是 $\infty$ 的共轭而 $\infty$ 是 $1$ 的共轭.

\begin{theorem}[H\"older 与 Minkowski 不等式]\label{theorem:ch19-holder-minkowski}
令 $(X,\mathcal{M},\mu)$ 为测度空间, $1\leqslant p<\infty$, 而 $q$ 是 $p$ 的共轭. 若 $f$ 属于 $L^{p}(X,\mu)$ 而 $g$ 属于 $L^{q}(X,\mu)$, 则它们的积 $fg$ 属于 $L^{1}(X,\mu)$.
\begin{description}
\item[H\"older 不等式]
\begin{align*}
\int_{X}|f\cdot g|\,d\mu=\|f\cdot g\|_{1}\leqslant\|f\|_{p}\cdot\|g\|_{q}.
\end{align*}
此外, 若 $f\neq 0$, 函数 $f^{*}=\|f\|_{p}^{1-p}\cdot\operatorname{sgn}(f)\cdot|f|^{p-1}$ 属于 $L^{q}(X,\mu)$, 且
\begin{align}
\int_{X}ff^{*}\,d\mu=\|f\|_{p}\text{ 且 }\|f^{*}\|_{q}=1.\label{align:ch19-holder-equality}
\end{align}
\item[Minkowski 不等式] 对于 $1\leqslant p\leqslant\infty$ 与 $f,g\in L^{p}(X,\mu)$,
\begin{align*}
\|f+g\|_{p}\leqslant\|f\|_{p}+\|g\|_{p}.
\end{align*}
因此 $L^{p}(X,\mu)$ 是一个赋范线性空间.
\end{description}
\end{theorem}

\begin{proof}
若 $p=1$, 则 H\"older 不等式从积分的单调性与齐次性, 以及观察到 $\|g\|_{\infty}$ 是 $g$ 的本性上界得出. 等式\eqref{align:ch19-holder-equality}是显然的. 假设 $p>1$. Young 不等式断言对非负实数 $a$ 和 $b$,
\begin{align*}
ab\leqslant\frac{1}{p}\cdot a^{p}+\frac{1}{q}\cdot b^{q}.
\end{align*}
若 $\|f\|_{p}=0$ 或 $\|g\|_{q}=0$, 则 $fg=0$ a.e., 不等式平凡成立. 假设 $\|f\|_{p}>0$ 且 $\|g\|_{q}>0$. 在 Young 不等式中取 $a=|f(x)|/\|f\|_{p}$ 和 $b=|g(x)|/\|g\|_{q}$, 积分得到
\begin{align*}
\frac{1}{\|f\|_{p}\|g\|_{q}}\int_{X}|fg|\,d\mu\leqslant\frac{1}{p\|f\|_{p}^{p}}\int_{X}|f|^{p}\,d\mu+\frac{1}{q\|g\|_{q}^{q}}\int_{X}|g|^{q}\,d\mu=\frac{1}{p}+\frac{1}{q}=1.
\end{align*}
因此 H\"older 不等式成立. 对于 $f^{*}$ 的验证是直接的.

Minkowski 不等式的证明与实直线上的情形相同(见 7.2 节定理 4). 对 $p=1$, 它从 $|f+g|\leqslant|f|+|g|$ 的积分得到. 对 $p>1$, 用 H\"older 不等式估计 $\int|f+g|^{p}=\int|f+g|^{p-1}|f+g|$.
\end{proof}

\begin{theorem}[Riesz-Fischer]\label{theorem:ch19-riesz-fischer}
令 $(X,\mathcal{M},\mu)$ 为测度空间而 $1\leqslant p\leqslant\infty$. 则 $L^{p}(X,\mu)$ 是 Banach 空间.
\end{theorem}

\begin{proof}
$p=\infty$ 的完备性的证明留作练习. 对于 $1\leqslant p<\infty$, 令 $\{f_{n}\}$ 为 $L^{p}(X,\mu)$ 中的 Cauchy 序列. 选取子序列 $\{f_{n_{k}}\}$ 使得 $\|f_{n_{k+1}}-f_{n_{k}}\|_{p}<1/2^{k}$. 定义
\begin{align*}
g_{k}=|f_{n_{1}}|+\sum_{j=1}^{k}|f_{n_{j+1}}-f_{n_{j}}|.
\end{align*}
则 $\{g_{k}\}$ 是上升的非负函数序列. 根据 Minkowski 不等式,
\begin{align*}
\|g_{k}\|_{p}\leqslant\|f_{n_{1}}\|_{p}+\sum_{j=1}^{k}\|f_{n_{j+1}}-f_{n_{j}}\|_{p}\leqslant\|f_{n_{1}}\|_{p}+1.
\end{align*}
因此 $\int g_{k}^{p}\,d\mu\leqslant(\|f_{n_{1}}\|_{p}+1)^{2}$. 根据单调收敛定理, $g=\lim g_{k}$ 满足 $\int g^{p}\,d\mu<\infty$. 因此 $g<\infty$ a.e., 从而级数
\begin{align*}
f_{n_{1}}+\sum_{j=1}^{\infty}(f_{n_{j+1}}-f_{n_{j}})
\end{align*}
在 $X$ 上 a.e. 绝对收敛. 定义 $f(x)$ 为该级数的和(若发散则为 $0$). 则 $\{f_{n_{k}}\}$ 在 $X$ 上 a.e. 逐点收敛于 $f$.

现在证明在 $L^{p}$ 中 $f_{n}\to f$. 令 $\varepsilon>0$. 选取 $N$ 使得对 $m,n\geqslant N$, $\|f_{m}-f_{n}\|_{p}<\varepsilon$. 对固定的 $k$ 使得 $n_{k}\geqslant N$, 根据 Fatou 引理,
\begin{align*}
\int|f-f_{n_{k}}|^{p}\,d\mu=\int\lim_{j\to\infty}|f_{n_{j}}-f_{n_{k}}|^{p}\,d\mu\leqslant\liminf_{j\to\infty}\int|f_{n_{j}}-f_{n_{k}}|^{p}\,d\mu\leqslant\varepsilon^{p}.
\end{align*}
因此 $f-f_{n_{k}}\in L^{p}$, 从而 $f\in L^{p}$, 且 $\|f-f_{n_{k}}\|_{p}\leqslant\varepsilon$. 由于 $\{f_{n}\}$ 是 Cauchy 的, 对 $n\geqslant N$,
\begin{align*}
\|f_{n}-f\|_{p}\leqslant\|f_{n}-f_{n_{k}}\|_{p}+\|f_{n_{k}}-f\|_{p}\leqslant 2\varepsilon.
\end{align*}
因此在 $L^{p}$ 中 $f_{n}\to f$.
\end{proof}

\begin{exercise}\label{exercise:ch19-01}
证明 $L^{\infty}(X,\mu)$ 是 Banach 空间.
\end{exercise}

\begin{exercise}\label{exercise:ch19-02}
令 $(X,\mathcal{M},\mu)$ 为测度空间, $1\leqslant p<\infty$, 而 $q$ 是 $p$ 的共轭. 证明简单函数在 $L^{p}(X,\mu)$ 中稠密.
\end{exercise}

\begin{exercise}\label{exercise:ch19-03}
令 $1\leqslant p<\infty$. 证明 H\"older 不等式中的等号成立当且仅当存在常数 $\alpha,\beta$, 使得 $\alpha|f|^{p}=\beta|g|^{q}$ a.e.
\end{exercise}

\begin{exercise}\label{exercise:ch19-04}
令 $\mu(X)<\infty$ 且 $1\leqslant p_{1}<p_{2}<\infty$. 证明 $L^{p_{2}}(X,\mu)\subseteq L^{p_{1}}(X,\mu)$. 给出一个反例说明当 $\mu(X)=\infty$ 时该包含关系不成立.
\end{exercise}

\begin{exercise}\label{exercise:ch19-05}
令 $\{f_{n}\}$ 为 $L^{p}(X,\mu)$ 中的序列, $1\leqslant p<\infty$, 使得在 $L^{p}$ 中 $f_{n}\to f$. 证明存在子序列 $\{f_{n_{k}}\}$ 在 $X$ 上 a.e. 逐点收敛于 $f$.
\end{exercise}

\end{document}